% Evans 4.3.1, printed pp. 177--186; PDF pp. 194--203.
\Needspace{65mm}
\section{Transform Methods}\label{sec:4-3}
\Needspace{35mm}
\subsection{Fourier transform: normalization and Plancherel}
For $f\in L^1(\R^n)$ define
\[
 \begin{aligned}
 \widehat f(\xi)&=(2\pi)^{-n/2}\int e^{-ix\cdot\xi}f(x)\,dx,\\
 \mathcal F^{-1}h(x)&=(2\pi)^{-n/2}\int e^{ix\cdot\xi}h(\xi)\,d\xi .
 \end{aligned}
\]
The integrals are absolutely convergent.
For Schwartz functions integration by parts and Fubini yield
\[
 \widehat{\partial_j f}=i\xi_j\widehat f,\qquad
 \widehat{f*h}=(2\pi)^{n/2}\widehat f\,\widehat h .
 \tag{F4}
\]
For the second identity, set $x=z+y$ in the double integral.
Both normalizations must be retained; $\mathcal F$ is unitary,
whereas ordinary convolution has no normalizing factor.
\Needspace{45mm}
\PDEAnnotation{卷积公式为何多一个常数}{每个 Fourier 变换自带 $(2\pi)^{-n/2}$，两个变换相乘就出现两份。卷积只做一次变换，所以要乘回 $(2\pi)^{n/2}$；后面的核常数都由此决定。}


\PDEBodyDiagram{4.2}{Fourier transformation turns differentiation into multiplication.}
{The square commutes with the unitary normalization used here. There is no $2\pi$ in the multiplier because the exponential is $e^{-ix\cdot\xi}$. The heat multiplier is small at large $|\xi|$, connecting this representation with the stability picture in Chapter 1.}
{ch04-fourier-square}

\begin{theorem}[Plancherel and inversion; source Theorems 1--2]
The Fourier transform extends uniquely to a unitary operator on $L^2$.
It preserves inner products, has the displayed inverse, and
$\widehat{D^\gamma f}=(i\xi)^\gamma\widehat f$ whenever
the weak derivative is in $L^2$.
\end{theorem}
\begin{expandedproof}[Gaussian regularization without circular inversion]
\paragraph{1. The Gaussian transform.}
For $a>0$, $I(b)=\int_{\R} e^{-ax^2}e^{-ibx}\,dx$.
Differentiate under the integral and integrate by parts to get
$I'(b)=-bI(b)/(2a)$.
Since $I(0)=\sqrt{\pi/a}$,
$I(b)=\sqrt{\pi/a}\,e^{-b^2/(4a)}$.
Products give the $n$-dimensional identity.
This real-variable proof does not yet require contour deformation.

\paragraph{2. The norm identity.}
For $f\in L^1\cap L^2$ and $\eta>0$, Fubini gives
\[
 \int e^{-\eta|\xi|^2}|\widehat f(\xi)|^2\,d\xi
 =\iint f(x)\overline{f(y)}G_\eta(x-y)\,dx\,dy,
\]
where $G_\eta(z)=(4\pi\eta)^{-n/2}e^{-|z|^2/(4\eta)}$.
Fubini is legitimate since $\|f\|_1^2<\infty$ and the frequency
Gaussian is integrable.
The right side tends to $\|f\|_2^2$ because $G_\eta*f\to f$ in $L^2$.
For completeness,
\[
 \|G_\eta*f-f\|_2
 \leq\int G_\eta(z)\|f(\cdot-z)-f\|_2\,dz\longrightarrow0:
\]
translations are continuous in $L^2$ by approximation with compact
continuous functions; split the last integral into a small ball and
its Gaussian tail. Monotone convergence on the left gives the norm identity.
Polarization gives the inner-product identity.

\paragraph{3. Inversion and extension.}
For Schwartz $f$, (F4) implies rapid decay of $\widehat f$.
Fubini and the Gaussian formula give
\[
 \mathcal F^{-1}(e^{-\eta|\xi|^2}\widehat f)=G_\eta*f.
\]
The right side tends to $f$; the left tends to
$\mathcal F^{-1}\widehat f$ by dominated convergence.
Thus inversion holds on Schwartz functions.
Approximate arbitrary $L^2$ data by Schwartz functions.
The norm identity makes their transforms Cauchy and makes the limit
independent of the approximation. Inversion extends by continuity,
so the range is all of $L^2$.
Finally integration by parts against Schwartz tests proves the derivative
identity in distributions; if both sides belong to $L^2$, it is an
a.e. equality of those representatives.
\end{expandedproof}
\Needspace{45mm}
\PDEAnnotation{$L^2$ 变换不必是逐点积分}{一般 $L^2$ 函数未必可积，其 Fourier 变换通过范数极限定义。反演和乘子公式首先是 $L^2$ 等式；若要逐点表示，还须增加可积性或正则性。}

\Needspace{35mm}
\subsection{Fourier multipliers and evolution equations}
\paragraph{Bessel potential.}
For $-\Delta u+u=f$, $f\in L^2$,
$\widehat u=(1+|\xi|^2)^{-1}\widehat f$.
This gives a unique $L^2$ distributional solution with
$\int(1+|\xi|^2)^2|\widehat u|^2<\infty$.
A rigorous kernel representation is
\[
 K(x)=\int_0^\infty e^{-s}G_s(x)\,ds,\qquad u=K*f .
\]
Tonelli gives $K\geq0$ and $\|K\|_1=1$;
its transform is $(2\pi)^{-n/2}/(1+|\xi|^2)$.
The kernel called $B$ in the source is $(2\pi)^{n/2}K$.
The singularity of $K$ at zero in some dimensions does not prevent
its $L^1$ integrability or convolution action on $L^2$.
\Needspace{45mm}
\PDEAnnotation{乘子有界不表示核处处有限}{$(1+|\xi|^2)^{-1}$ 给出有界的 $L^2$ 乘子。对应核可能在原点无界，但仍可积；不能把一个点的奇性误当作表示公式失效。}

\paragraph{Heat and Schr\"odinger equations.}
Spatial transformation of $u_t-\Delta u=0$ gives
$\widehat u_t+|\xi|^2\widehat u=0$, hence
$\widehat u=e^{-t|\xi|^2}\widehat g$.
Gaussian inversion recovers $u=G_t*g$ from Chapter~2.
For $iu_t+\Delta u=0$ the sign is instead
$\widehat u_t=-i|\xi|^2\widehat u$, so
\[
 \widehat u=e^{-it|\xi|^2}\widehat g,\qquad
 u(x,t)=(4\pi it)^{-n/2}\int e^{i|x-y|^2/(4t)}g(y)\,dy .
 \tag{SC4}
\]
For $t\ne0$ use
$(it)^{-n/2}=|t|^{-n/2}e^{-i\pi n\,\operatorname{sgn}(t)/4}$.
The kernel identity follows by first replacing $it$ with $\eta+it$,
$\eta>0$, and then letting $\eta\downarrow0$.
For $g\in L^1\cap L^2$, the kernels converge under the spatial integral
by domination, while the multipliers converge in $L^2$.
Thus both descriptions agree.
For $L^2$ data this defines a unitary group:
\[
 \|u(t)\|_2=\|g\|_2,\qquad
 \|u(t)-g\|_2^2
 =\int|e^{-it|\xi|^2}-1|^2|\widehat g|^2\longrightarrow0 .
\]
If $g,|x|^2g\in L^1$, differentiation of (SC4) on compact
space-time sets away from $t=0$ verifies the classical PDE.
For arbitrary $L^2$ data the statement is distributional.
\Needspace{45mm}
\PDEAnnotation{色散与耗散的范数表现不同}{热乘子模长小于一，会压低高频；Schr\"odinger 乘子模长为一，保持 $L^2$。后者的点态衰减来自振荡抵消，不来自 $L^2$ 质量损失。}

\paragraph{Wave and telegraph equations.}
For $u_{tt}-\Delta u=0$, $u(0)=g$, $u_t(0)=h$,
\[
 \widehat u(\xi,t)=\cos(t|\xi|)\widehat g
              +\frac{\sin(t|\xi|)}{|\xi|}\widehat h.
 \tag{WV4}
\]
At $\xi=0$ the quotient is defined as $t$.
It is bounded by $t$, so the formula has no genuine zero-frequency pole.
For $u_{tt}+2du_t-u_{xx}=0$ use $\widehat u=e^{-dt}v$.
Then $v''+(\xi^2-d^2)v=0$ and
\[
 \widehat u=e^{-dt}
 \begin{cases}
 \cosh(qt)\widehat g+\dfrac{\sinh(qt)}q(\widehat h+d\widehat g),
                 &|\xi|<d,\ q=\sqrt{d^2-\xi^2},\\
 \widehat g+t(\widehat h+d\widehat g),&|\xi|=d,\\
 \cos(qt)\widehat g+\dfrac{\sin(qt)}q(\widehat h+d\widehat g),
                 &|\xi|>d,\ q=\sqrt{\xi^2-d^2}.
 \end{cases}
\]
These are the overdamped, critical and oscillatory regimes.
The common limiting formula at $q=0$ is essential.

\begin{theorem}[Asymptotic energy equipartition; source \S4.3.1]
If $g\in H^1(\R^n)$ and $h\in L^2(\R^n)$, the wave solution has
\[
 \|Du(t)\|_2^2\longrightarrow E_0,\qquad
 \|u_t(t)\|_2^2\longrightarrow E_0,\qquad
 E_0=\tfrac12(\|Dg\|_2^2+\|h\|_2^2).
\]
\end{theorem}
\begin{expandedproof}
Write $A=|\xi|\widehat g$ and $B=\widehat h$ in $L^2$.
Plancherel and (WV4) give
\[
 \|Du\|_2^2
 =\int|A|^2\cos^2(t|\xi|)+|B|^2\sin^2(t|\xi|)
       +2\operatorname{Re}(A\overline B)\sin(t|\xi|)\cos(t|\xi|).
\]
All three coefficient functions are $L^1$, the cross term by
Cauchy--Schwarz.
For any $F\in L^1(\R^n)$, polar coordinates give
$\int e^{2it|\xi|}F(\xi)\,d\xi=\int_0^\infty e^{2itr}H(r)\,dr$
with $\|H\|_1\leq\|F\|_1$.
The latter tends to zero: approximate $H$ in $L^1$ by smooth functions
supported inside $(0,\infty)$ and integrate those by parts.
Now use the double-angle identities to obtain the first limit.
The analogous formula for $u_t$ gives energy conservation and
the second limit.
This also handles $n=1$ without discarding a possible endpoint term
in a radial integration by parts.
\end{expandedproof}
% Supplement to Evans 4.3: Radon and Laplace transforms, printed pp. 186--195.
\Needspace{35mm}
\subsection{Radon transform and plane-wave decomposition}\label{sec:4-radon}
For $f\in C_c^\infty(\R^n)$, set
\[
 Rf(s,\omega)=\int_{x\cdot\omega=s}f(x)\,dS_x,\qquad \omega\in S^{n-1}.
\]
Coordinates $x=s\omega+\sum y_jb_j$, where $b_j$ are an orthonormal
basis of $\omega^\perp$, make the surface measure $dy$.
They immediately give
\[
 Rf(-s,-\omega)=Rf(s,\omega),\quad
 R(\partial_jf)=\omega_j\partial_sRf,\quad R(\Delta f)=\partial_s^2Rf .
\]
For the middle identity decompose $\partial_j$ into a normal and
tangential derivative. Tangential integrals vanish by compact support.
Induction gives $R(D^\gamma f)=\omega^\gamma\partial_s^{|\gamma|}Rf$.
If $\operatorname{supp}f\subset B_R$, then $Rf=0$ for $|s|\geq R$.


\PDEBodyDiagram{4.3}{Radon data integrate along level hyperplanes.}
{In dimension two the integration sets are lines perpendicular to $\omega$. The scalar $s$ is signed distance in the normal direction, not distance traveled along a line. The blue line is one member of the family; higher dimensions replace it by an $(n-1)$-plane.}
{ch04-radon-lines}

\begin{theorem}[Slice theorem and inversion; source Theorems 4--5]
With the unnormalized one-dimensional transform
$B(r,\omega)=\int_{\R} e^{-irs}Rf(s,\omega)\,ds$,
\[
 B(r,\omega)=(2\pi)^{n/2}\widehat f(r\omega).
\]
For every dimension,
\[
 f(x)=\frac1{2(2\pi)^n}
 \int_{S^{n-1}}\int_{\R} B(r,\omega)|r|^{n-1}e^{irx\cdot\omega}\,dr\,dS_\omega .
\]
If $n=2k+1$ is odd, this becomes the local formula
\[
 f(x)=C_n\int_{S^{n-1}}
       \partial_s^{2k}Rf(x\cdot\omega,\omega)\,dS_\omega,\qquad
 C_n=\frac{(-1)^k}{2(2\pi)^{2k}} .
 \tag{RI4}
\]
\end{theorem}
\begin{expandedproof}
Fubini in the orthonormal coordinates gives
\[
 B(r,\omega)=\int f(x)e^{-ir\omega\cdot x}\,dx .
\]
In Fourier inversion, polar coordinates with signed $r$ double-count
each nonzero frequency, giving the factor $1/2$ and $|r|^{n-1}$.
All integrals converge since $\widehat f$ is rapidly decreasing.
For $n-1=2k$, one-dimensional inversion gives
\[
 \partial_s^{2k}Rf(s,\omega)
 =\frac{(-1)^k}{2\pi}\int_{\R} r^{2k}B(r,\omega)e^{irs}\,dr .
\]
Substituting this in the general inversion formula proves (RI4).
\end{expandedproof}
\Needspace{45mm}
\PDEAnnotation{奇偶维差别藏在绝对值里}{带符号的极径在一般维数必须用 $|r|^{n-1}$；只有奇维时才等于多项式 $r^{2k}$，从而变成有限阶局部导数。偶维的反演通常带非局部算子。}

For odd $n$, define $\rho_f=C_n\partial_s^{n-1}Rf$.
Then the plane-wave integral
\[
 u(x,t)=\int_{S^{n-1}}\rho_g(x\cdot\omega-t,\omega)\,dS_\omega
\]
solves $u_{tt}-\Delta u=0$, $u(0)=g$, $u_t(0)=0$.
The PDE holds because $|\omega|=1$; the first initial condition is
(RI4). The integrand for $u_t(0)$ is odd under $\omega\mapsto-\omega$,
since an odd $s$ derivative reverses the symmetry of $Rg$.

For the constant-coefficient system
$\mathbf u_t+\sum B_j\mathbf u_{x_j}=0$,
suppose $B(\omega)=\sum\omega_jB_j$ is diagonalizable with real
eigenvalues $\lambda_l(\omega)$.
In Radon coordinates the diagonal components satisfy one-dimensional
transport and equal
\[
 v_l(s,\omega,t)=h_l(s-\lambda_l(\omega)t,\omega).
\]
After applying (RI4), use
$\rho_l(s,\omega)=C_n\partial_s^{n-1}h_l(s,\omega)$ exactly once:
\[
 \mathbf u(x,t)=\int_{S^{n-1}}\sum_l
 a_l(\omega)\rho_l(x\cdot\omega-\lambda_l(\omega)t,\omega)\,dS_\omega .
\]
Here $a_l$ are the diagonalizing columns; the shift is not inserted
again in the definition of $\rho_l$.
For compact data in $B_R$, if
$\min_{\omega,l}|\lambda_l(\omega)|=\sigma>0$,
then $\mathbf u=0$ for $|x|<\sigma t-R$.
Indeed each shifted argument has magnitude at least $\sigma t-|x|>R$.
This is a conditional odd-dimensional Huygens statement; the
diagonalization and nonzero-speed hypotheses cannot be removed.
\Needspace{45mm}
\PDEAnnotation{平面支集条件怎样变成内部空区}{若 $|x|<\sigma t-R$，所有方向的平面参数都移出初始支集区间，整个反演积分为零。零特征速度会留下不移动的分量，因此必须单独排除。}

\Needspace{35mm}
\subsection{Laplace transform, resolvents and the heat--wave link}
If $v$ and $v_t$ have suitable exponential bounds, integration by parts gives
\[
 \mathcal L(v_t)(s)=s\mathcal Lv(s)-v(0),\qquad
 \mathcal Lv(s)=\int_0^\infty e^{-st}v(t)\,dt .
\]
Thus for $v_t-\Delta v=0$, $v(0)=f$,
\[
 (s-\Delta)\mathcal Lv=f .
\]
Use the same spatial boundary condition before and after transformation;
for example a zero Dirichlet condition stays zero.
In the whole space this recovers
$(s-\Delta)^{-1}f=\int_0^\infty e^{-st}(G_t*f)\,dt$.
For $L^2$ data the Bochner integral is justified by
$\|G_t*f\|_2\leq\|f\|_2$.
\Needspace{45mm}
\PDEAnnotation{初值正是 Laplace 分部积分的边界项}{时间导数并非只替换成乘以 $s$。零时刻留下 $-v(0)$，它成为椭圆 resolvent 方程的右端；漏掉这一项就把非零初值问题变成了齐次问题。}

\begin{theorem}[Gaussian averaging of the wave group; source Example 9]
Let $g$ be smooth compactly supported, and let $U(x,s)$ be the wave
solution with displacement $g$ and zero initial velocity, extended evenly
in $s$. Then
\[
 V(x,t)=\frac1{\sqrt{4\pi t}}\int_{\R} e^{-s^2/(4t)}U(x,s)\,ds
\]
is the heat solution with datum $g$.
\end{theorem}
\begin{expandedproof}
For smooth data the derivatives of $U$ grow at most polynomially,
as can also be checked from its Fourier representation.
Gaussian weights justify differentiation and kill the boundary terms.
Writing $q_t(s)=(4\pi t)^{-1/2}e^{-s^2/(4t)}$, observe
$\partial_tq_t=\partial_{ss}q_t$. Hence
\[
 V_t=\int q_{ss}U=\int qU_{ss}=\int q\Delta U=\Delta V.
\]
The Gaussian approximate identity in $s$ gives $V(0)=g$.
Alternatively the frequency multiplier is
$\int q_t(s)\cos(s|\xi|)\,ds=e^{-t|\xi|^2}$,
which proves the same statement in $L^2$.
\end{expandedproof}

For $n=2k+1\geq3$, let $M_g(x,r)$ be the spherical average of $g$.
Equating the two Gaussian representations gives, with
$\lambda=1/(4t)$,
\[
 \int_0^\infty e^{-\lambda s^2}U(x,s)\,ds
 =\frac{|S^{n-1}|}{2\pi^k}\lambda^k
       \int_0^\infty e^{-\lambda r^2}r^{2k}M_g(x,r)\,dr .
\]
Use $\lambda e^{-\lambda r^2}=-(2r)^{-1}\partial_re^{-\lambda r^2}$
and integrate by parts $k$ times. The result is
\[
 \lambda^k\int_0^\infty e^{-\lambda r^2}r^{2k}M_g\,dr
 =2^{-k}\int_0^\infty e^{-\lambda r^2}
       r(r^{-1}\partial_r)^k(r^{2k-1}M_g)\,dr .
\]
There are no boundary terms at infinity by compact support in $r$.
At zero, $M_g$ is smooth and even; before the last differentiation
the relevant odd power has not fallen below one, so the boundary
terms are zero.
After $s^2=\tau$, uniqueness of the Laplace transform implies
\[
 U(x,t)=\frac1{(n-2)!!}\,
 \partial_t(t^{-1}\partial_t)^{(n-3)/2}
                  (t^{n-2}M_g(x,t)).
\]
Here $|S^{n-1}|/(2^{k+1}\pi^k)=1/(n-2)!!$ follows from the
Gaussian sphere area formula.
The one-dimensional case is $U=(g(x+t)+g(x-t))/2$.
One can justify Laplace uniqueness without a new PDE theorem:
after multiplying by $e^{-a\tau}$, substitute $z=e^{-\tau}$;
vanishing at all $s=a+j$ gives zero moments on $[0,1]$.
Polynomial density in continuous functions makes the finite signed
measure zero. This applies to the continuous, exponentially integrable
functions above.

